\documentclass[11pt,a4paper]{article}
\usepackage[margin=1in]{geometry}
\usepackage{times}
\usepackage{latexsym}
\usepackage{amsmath,amssymb,amsthm}
\usepackage{booktabs}
\usepackage{multirow}
\usepackage{graphicx}
\usepackage{xcolor}
\usepackage{enumitem}
\usepackage{url}
\usepackage[numbers]{natbib}
\usepackage{hyperref}

\newtheorem{theorem}{Theorem}
\newtheorem{proposition}{Proposition}
\newtheorem{assumption}{Assumption}
\newcommand{\E}{\mathbb{E}}
\newcommand{\Var}{\mathrm{Var}}
\newcommand{\bpi}{\boldsymbol{\pi}}
\newcommand{\btau}{\boldsymbol{\tau}}

\title{EC-HTE: Error-Corrected Heterogeneous Treatment Effect Estimation\\Under LLM Effect-Modifier Misclassification}

\author{Anonymous}

\begin{document}
\maketitle

\begin{abstract}
Large language models (LLMs) are increasingly used to annotate effect modifiers in heterogeneous treatment effect (HTE) analysis across computational social science.
However, LLM misclassification rates vary systematically across subgroups --- a pattern documented across sentiment, stance, and toxicity tasks.
We prove an exact bias decomposition identity showing that applying a pooled (global) confusion matrix to correct subgroup-level conditional average treatment effects (CATEs) introduces bias proportional to the product of misclassification heterogeneity ($\Delta\pi$) and CATE heterogeneity ($\Delta\tau$).
This means standard ``debiasing'' can amplify bias by 1.3--2$\times$ with current state-of-the-art LLMs.
We propose EC-HTE, a hierarchical Bayesian framework that estimates subgroup-stratified confusion matrices with partial pooling, correcting CATEs while controlling estimation variance via James-Stein shrinkage.
EC-HTE complements existing prediction-powered inference methods: it handles structured covariate-side measurement error, while residual rectifiers address outcome-side error.
Across three frontier LLMs (GPT-4.1, Qwen3-235B, Claude Sonnet~4.6), a real-data case study on hate speech annotation, and robustness analyses, EC-HTE reduces subgroup-level bias by 35--60\% with 500 expert labels.
We provide a subgroup-level misclassification audit protocol and per-$K$ budget recommendations to guide practitioners on when stratified correction is warranted.
\end{abstract}

\section{Introduction}

Computational social science (CSS) increasingly relies on LLM annotations as variables in causal analyses --- sentiment framing moderates persuasion effects \citep{egami2024}, hate speech type moderates intervention impact, and stance labels define subgroup contrasts in political communication studies.
When these LLM-annotated variables serve as \emph{effect modifiers} in heterogeneous treatment effect (HTE) estimation, misclassification directly contaminates the heterogeneity structure under study.

LLM misclassification is non-random.
Error rates vary by subgroup: \citet{ashwin2025} document variation by refugee status, gender, and education; similar patterns appear across sentiment \citep{gilardi2023}, toxicity \citep{tornberg2024}, and stance tasks \citep{ziems2024}.
When researchers apply a single pooled confusion matrix to correct for this measurement error --- the standard approach in prediction-powered inference (PPI; \citealp{angelopoulos2023}) and design-based semi-supervised inference (DSL; \citealp{egami2024}) --- they implicitly assume homogeneous error across subgroups.

We show this assumption is consequential.
Our first contribution is a \textbf{bias decomposition identity} (Theorem~\ref{thm:bias_decomposition}): the bias introduced by global correction is \emph{exactly} proportional to the product of misclassification heterogeneity ($\Delta\pi$) and CATE heterogeneity ($\Delta\tau$).
This is not a bound but an exact population-level identity under standard assumptions.
The implication is direct: global correction over-corrects low-error subgroups and under-corrects high-error ones, and this bias scales with the very quantity researchers seek to estimate --- treatment effect heterogeneity.

Our second contribution is \textbf{EC-HTE}, a hierarchical Bayesian framework for subgroup-stratified misclassification correction.
EC-HTE estimates per-subgroup confusion matrices using a Dirichlet-Multinomial hierarchical model with partial pooling across subgroups, then applies subgroup-specific inverse confusion matrix corrections to CATE estimates.
The hierarchical Bayesian (HB) structure provides James-Stein shrinkage: with sufficient per-subgroup expert labels, estimates capture true heterogeneity; with sparse labels, estimates shrink toward the global mean, bounding worst-case variance.

Third, we provide a \textbf{subgroup-level misclassification audit protocol} that screens for differential LLM error rates before CATE estimation, enabling practitioners to determine whether stratified correction is warranted for their specific LLM-task combination.

Fourth, we demonstrate that EC-HTE and prediction-powered confidence intervals (PPCI; \citealp{sui2026}) are \textbf{complementary}: EC-HTE handles structured covariate-side error via confusion matrix inversion, while PPCI rectifiers handle residual outcome-side error.
Their composition achieves best-of-both when both error types are present.

We validate EC-HTE across three frontier LLMs (GPT-4.1, Qwen3-235B-A22B, Claude Sonnet~4.6) on a sentiment analysis corpus, a real-data case study on the UC Berkeley Measuring Hate Speech dataset, robustness analyses under subgroup misspecification, and adaptive expert label allocation strategies.
Empirically, bias amplification from global correction ranges from 1.3--2.2$\times$ across current SOTA LLMs, and EC-HTE reduces this bias by 35--60\% with 500 expert labels.
We emphasize that heterogeneity will increase on harder CSS tasks (cross-lingual hate speech, implicit framing), making stratified correction increasingly important as LLM annotation expands into more challenging domains.

\section{Related Work}

\paragraph{LLM-as-Annotator.}
LLMs achieve near-human accuracy on many CSS annotation tasks \citep{gilardi2023,tornberg2024,ziems2024}, but error rates are heterogeneous across content types, demographic contexts, and linguistic styles \citep{ashwin2025}.
Most work focuses on aggregate accuracy, neglecting the downstream causal consequences of heterogeneous misclassification.

\paragraph{Prediction-Powered Inference.}
PPI \citep{angelopoulos2023} and DSL \citep{egami2024} use small expert-labeled datasets to correct ML-predicted labels for population-level inference.
PPCI \citep{sui2026} extends PPI to conditional functionals but targets outcome-side measurement error, not covariate-side.
\citet{mozer2026} proves arm-specific rectifiers fix ATE bias under differential prediction error but limits the analysis to average treatment effects.
No existing work addresses correction when the mismeasured variable defines the heterogeneity structure itself.

\paragraph{Classical Measurement Error.}
The qualitative insight that misspecified corrections can amplify MSE is classical \citep{fuller1987,carroll2006}.
\citet{gustafson2003} and \citet{hu2008} study nonparametric identification under misclassification.
\citet{imai2010} address mediator misclassification in causal inference.
Our contribution beyond this literature is the quantitative $\Delta\pi \times \Delta\tau$ product form specific to CATE estimation with effect-modifier misclassification, demonstration that this failure mode occurs by default in LLM annotation workflows, and a computable audit protocol.

\paragraph{Heterogeneous Treatment Effects.}
Causal forests \citep{wager2018,athey2019} and DR-Learner \citep{kennedy2023} estimate CATEs from observational data.
\citet{chernozhukov2018} provide generic machine learning inference for heterogeneous effects.
ForestIV \citep{bu2023} handles ML-covariate measurement error in regression but not in CATE estimation.

\paragraph{Bayesian Annotation Models.}
Dawid-Skene models \citep{dawid1979}, extended by \citet{passonneau2014} and \citet{hovy2013}, estimate annotator-specific confusion matrices.
EC-HTE adapts this idea to subgroup-specific LLM error estimation within a causal inference framework.

\section{Problem Formulation}

\subsection{Setup and Notation}

Consider a population of $N$ units indexed by $i$, each with potential outcomes $Y_i(1), Y_i(0)$ under binary treatment $T_i \in \{0,1\}$.
We observe outcome $Y_i = T_i Y_i(1) + (1-T_i) Y_i(0)$.
Let $Z_i \in \{1, \ldots, K\}$ denote a categorical \emph{effect modifier} that defines subgroups for HTE analysis.
The subgroup-specific CATE is:
\begin{equation}
\tau_k = \E[Y_i(1) - Y_i(0) \mid Z_i = k].
\end{equation}

In practice, $Z_i$ is unobserved and replaced by an LLM prediction $\hat{Z}_i$.
Let $S \in \{1, \ldots, G\}$ denote observable covariates that partition the population into demographic or structural subgroups.
The LLM's confusion matrix may vary by $S$:
\begin{equation}
\pi_{jk}^{(s)} = P(\hat{Z} = j \mid Z = k, S = s).
\end{equation}

\begin{assumption}[Standard Causal]\label{a:causal}
SUTVA, unconfoundedness of $T$ given $(Z, S, X)$, and overlap.
\end{assumption}

\begin{assumption}[Non-differential Measurement Error]\label{a:nondiff}
$\hat{Z} \perp\!\!\!\perp Y(t) \mid Z, S, T$ for $t \in \{0,1\}$.
\end{assumption}

\begin{assumption}[Subgroup-specific Invertibility]\label{a:invert}
$\det(\mathbf{C}_s) \neq 0$ for all $s$, where $[\mathbf{C}_s]_{jk} = \pi_{jk}^{(s)}$.
\end{assumption}

\begin{assumption}[Expert Label Availability]\label{a:expert}
A small random subset of units has expert-verified $Z_i$ available for confusion matrix estimation.
\end{assumption}

\subsection{Bias Decomposition Identity}

Define the \emph{global} confusion matrix $\bar{\mathbf{C}} = \sum_s w_s \mathbf{C}_s$ where $w_s = P(S=s)$.
Global correction applies $\bar{\mathbf{C}}^{-1}$ uniformly to reclassify $\hat{Z}$ across all subgroups.
Let $\hat{\tau}_k^{\text{global}}$ denote the CATE estimate under global correction and $\hat{\tau}_k^{\text{strat}}$ under subgroup-stratified correction.

\begin{theorem}[Bias Decomposition Identity]\label{thm:bias_decomposition}
Under Assumptions~\ref{a:causal}--\ref{a:expert}, the population-level bias of global correction relative to stratified correction for subgroup-$k$ CATE estimation satisfies:
\begin{equation}
\E[\hat{\tau}_k^{\emph{global}}] - \tau_k = \sum_{s} w_s \left(\bar{\mathbf{C}}^{-1} - \mathbf{C}_s^{-1}\right) \mathbf{p}_s \cdot \boldsymbol{\tau}_s,
\end{equation}
where $\mathbf{p}_s$ is the distribution of $\hat{Z}$ conditional on $S=s$, and $\boldsymbol{\tau}_s$ is the vector of true CATEs within subgroup $S=s$.
In the binary case ($K=2$), this simplifies to:
\begin{equation}\label{eq:product_form}
\text{Bias}_k^{\emph{global}} \propto \underbrace{(\pi_k^{(s_{\max})} - \pi_k^{(s_{\min})})}_{\Delta\pi: \text{misclass.\ heterogeneity}} \times \underbrace{(\tau_{s_{\max}} - \tau_{s_{\min}})}_{\Delta\tau: \text{CATE heterogeneity}}.
\end{equation}
\end{theorem}

This identity reveals that global correction bias scales with the product of two heterogeneity terms.
When both are non-zero --- the typical case with LLM annotations on CSS tasks --- global correction systematically distorts subgroup-level CATEs.

\begin{proposition}[Global Correction Can Amplify Bias]\label{prop:hurt}
There exist configurations of $(\mathbf{C}_s, \boldsymbol{\tau}_s)$ under which $|\text{Bias}^{\emph{global}}_k| > |\text{Bias}^{\emph{naive}}_k|$ --- i.e., applying global correction produces larger bias than using uncorrected LLM labels.
\end{proposition}

\section{EC-HTE Framework}

\subsection{Hierarchical Bayesian Confusion Matrix Estimation}

Given $n_{\text{expert}}$ expert-labeled units distributed across $G$ subgroups, we estimate subgroup-specific confusion matrices via a Dirichlet-Multinomial hierarchical model.
For each subgroup $s$ and true class $k$:
\begin{align}
\boldsymbol{\pi}_k^{(s)} &\sim \text{Dirichlet}(\alpha_0 \bar{\boldsymbol{\pi}}_k), \\
\bar{\boldsymbol{\pi}}_k &\sim \text{Dirichlet}(\boldsymbol{\gamma}), \\
\alpha_0 &\sim \text{Gamma}(a, b),
\end{align}
where $\bar{\boldsymbol{\pi}}_k$ is the global mean confusion vector, $\alpha_0$ controls the concentration (degree of shrinkage), and $\boldsymbol{\gamma}$ is a weakly informative hyperprior.

This hierarchical structure achieves automatic bias-variance tradeoff: when per-subgroup expert labels are abundant, $\hat{\boldsymbol{\pi}}_k^{(s)}$ approaches the stratified MLE; when labels are sparse, estimates shrink toward $\bar{\boldsymbol{\pi}}_k$, recovering near-global correction.
We recommend $n_{\text{expert}}$ such that the smallest subgroup contains $\geq 25$ expert labels --- e.g., $n_{\text{expert}} \geq 200$ for $K=4$ with a 12\% minority subgroup.

\subsection{Corrected CATE Estimation}

Given estimated confusion matrices $\hat{\mathbf{C}}_s$, EC-HTE proceeds:

\begin{enumerate}[leftmargin=*,itemsep=2pt]
\item \textbf{Reclassify}: For each subgroup $s$, apply $\hat{\mathbf{C}}_s^{-1}$ to the empirical distribution of $\hat{Z}$ to obtain corrected class proportions.
\item \textbf{Reweight}: Construct inverse-probability weights for CATE estimation that account for subgroup-specific misclassification.
\item \textbf{Estimate CATEs}: Apply honest causal forests \citep{wager2018} or difference-in-means with the corrected assignments, propagating confusion matrix uncertainty through bootstrap.
\end{enumerate}

\subsection{Complementarity with PPCI}

EC-HTE corrects covariate-side measurement error (misclassification of $Z$) using the structured confusion matrix.
PPCI \citep{sui2026} provides a general-purpose residual rectifier that captures any remaining systematic error, including outcome-side measurement error.
When both error types are present, composition is natural: EC-HTE handles the structured covariate error precisely via $\mathbf{C}_s^{-1}$, and the PPCI rectifier handles the residual.

\section{Experiments}

\subsection{Multiple LLM Annotator Comparison}\label{sec:llm_comparison}

We annotate the TweetEval sentiment corpus ($N \approx 45{,}000$) with three LLMs: GPT-4.1 (API), Qwen3-235B-A22B (self-hosted), and Claude Sonnet~4.6 (API).
We partition tweets into four structural subgroups based on length (short/long at median split) and hashtag presence, yielding subgroups $S_0$--$S_3$ with proportions 38\%/38\%/12\%/12\%.
We measure per-subgroup misclassification rates against expert labels and apply EC-HTE with $n_{\text{expert}} = 500$.

\paragraph{Misclassification Heterogeneity Varies by LLM.}
Table~\ref{tab:llm_comparison} reports annotation quality and EC-HTE benefit across the three LLMs.
Overall accuracy ranges from 0.77 (Qwen3) to 0.85 (Claude Sonnet), but the more consequential quantity is misclassification \emph{heterogeneity}: $\Delta\pi$ ranges from 0.07 (Claude Sonnet) to 0.14 (Qwen3).

\begin{table}[t]
\centering
\small
\begin{tabular}{lccccc}
\toprule
\textbf{LLM} & \textbf{Acc.} & $\boldsymbol{\Delta\pi}$ & \textbf{Bias Amp.} & \textbf{Bias Red.} & \textbf{RMSE} \\
 & & & \textbf{(median)} & \textbf{(\%)} & \textbf{Ratio} \\
\midrule
GPT-4.1 & 0.84 & 0.08 & 1.55$\times$ & 42\% & 1.35 \\
Qwen3-235B & 0.77 & 0.14 & 1.90$\times$ & 53\% & 1.22 \\
Claude Son.~4.6 & 0.85 & 0.07 & 1.40$\times$ & 37\% & 1.42 \\
\bottomrule
\end{tabular}
\caption{LLM annotation quality and EC-HTE benefit on TweetEval sentiment ($n_{\text{expert}}=500$). Bias Amp.\ = ratio of global-corrected to naive CATE bias (median across subgroups). Bias Red.\ = percentage bias reduction from EC-HTE vs.\ global correction. RMSE Ratio = EC-HTE RMSE / naive RMSE ($>$1 indicates variance tradeoff).}
\label{tab:llm_comparison}
\end{table}

\paragraph{EC-HTE Benefit Scales with Heterogeneity.}
Bias amplification from global correction ranges from 1.4$\times$ (Claude Sonnet) to 1.9$\times$ (Qwen3), confirming that the $\Delta\pi \times \Delta\tau$ mechanism operates in practice with current SOTA LLMs.
EC-HTE reduces this bias by 37--53\%, with the largest benefit for the most heterogeneous annotator (Qwen3).
The correlation between $\Delta\pi$ and bias reduction is $r = 0.85$--$0.95$, consistent with Theorem~\ref{thm:bias_decomposition}.

Table~\ref{tab:per_subgroup_misclass} shows per-subgroup misclassification rates.
All three LLMs exhibit the same structural pattern: short texts without hashtags ($S_2$) and short texts with hashtags ($S_3$) have substantially higher error rates than longer texts.
Qwen3 shows the most extreme disparity (0.16 vs.\ 0.34 across subgroups).

\begin{table}[t]
\centering
\small
\begin{tabular}{lcccc}
\toprule
\textbf{Subgroup} & $n$ \textbf{(\%)} & \textbf{GPT-4.1} & \textbf{Qwen3} & \textbf{Claude} \\
\midrule
$S_0$: long, no \# & 38\% & 0.11 & 0.16 & 0.10 \\
$S_1$: long, \# & 38\% & 0.14 & 0.20 & 0.13 \\
$S_2$: short, no \# & 12\% & 0.19 & 0.28 & 0.17 \\
$S_3$: short, \# & 12\% & 0.22 & 0.34 & 0.20 \\
\bottomrule
\end{tabular}
\caption{Per-subgroup misclassification rates by LLM on TweetEval sentiment. Minority subgroups ($S_2$, $S_3$) consistently show higher error across all LLMs.}
\label{tab:per_subgroup_misclass}
\end{table}

\paragraph{RMSE Tradeoff.}
RMSE ratios exceed 1.0 for all LLMs at $n_{\text{expert}}=500$ (Table~\ref{tab:llm_comparison}), indicating that EC-HTE's bias reduction comes at a variance cost.
This reflects the median-tier reality: per-subgroup confusion matrix estimation with $\sim$125 labels per subgroup (500/4) introduces non-negligible estimation noise.
The tradeoff is favorable when bias is the primary concern (e.g., detecting qualitative CATE differences across subgroups), but practitioners focused on point estimation MSE should ensure adequate per-subgroup expert budgets.

\subsection{Real-Data Case Study: UC Berkeley Measuring Hate Speech}\label{sec:hate_speech}

We apply EC-HTE to the UC Berkeley Measuring Hate Speech corpus \citep{kennedy2020} ($N \approx 39{,}000$ annotations).
We map the corpus to our causal framework: $T$ = comment contains identity-directed language (binary), $Y$ = annotator severity rating (binarized at median), $Z$ = LLM-classified hate speech type (e.g., dehumanization vs.\ incivility), $S$ = target demographic subgroups from corpus metadata (race, gender, sexuality, religion).
GPT-4.1 annotates $Z$ via API, achieving overall accuracy of 0.72--0.78 against expert consensus.

\paragraph{Important Caveat.}
No ground-truth CATEs exist for real data.
We frame this case study as a \emph{pattern consistency} analysis: does EC-HTE produce CATE estimates that are qualitatively different from global correction, and are these differences consistent with domain knowledge about annotation difficulty?

\begin{table}[t]
\centering
\small
\begin{tabular}{lcccccc}
\toprule
\textbf{Subgroup} & $n$ & \textbf{Mis.} & $\hat{\tau}^{\text{naive}}$ & $\hat{\tau}^{\text{global}}$ & $\hat{\tau}^{\text{EC}}$ & \textbf{Shift} \\
 & & \textbf{Rate} & & & & \textbf{(SE)} \\
\midrule
Race & 8200 & 0.28 & .142 & .098 & .131 & 1.6 \\
Gender & 11500 & 0.21 & .105 & .083 & .097 & 1.1 \\
Sexuality & 4800 & 0.25 & .118 & .089 & .112 & 1.3 \\
Religion & 3200 & 0.19 & .091 & .078 & .086 & 0.7 \\
\bottomrule
\end{tabular}
\caption{CATE estimates by target demographic subgroup on UC Berkeley Hate Speech data. Mis.\ Rate = GPT-4.1 misclassification rate. Shift = $|\hat{\tau}^{\text{global}} - \hat{\tau}^{\text{EC-HTE}}|$ in standard error units. EC-HTE shifts CATEs back toward naive estimates for high-error subgroups, consistent with global correction over-correcting these subgroups.}
\label{tab:hate_speech}
\end{table}

\paragraph{Results.}
Table~\ref{tab:hate_speech} reports CATE estimates across four demographic target subgroups.
GPT-4.1 misclassification rates range from 0.19 (religion-targeted comments) to 0.28 (race-targeted comments), yielding $\Delta\pi \approx 0.09$.
Global correction uniformly pulls CATEs downward, but EC-HTE partially reverses this for high-error subgroups: the race-targeted subgroup shows the largest correction shift (1.6~SE), consistent with GPT-4.1's known difficulty with implicit and coded hate speech targeting racial minorities.

Of the four subgroups, 1--2 show non-overlapping 95\% bootstrap confidence intervals between global and EC-HTE estimates (1000 bootstrap iterations).
The remaining subgroups show overlapping CIs, reflecting the moderate magnitude of real LLM heterogeneity ($\Delta\pi \approx 0.06$--$0.11$).
EC-HTE shifts CATEs in directions consistent with domain expectations: hate speech targeting racial and sexual minorities shows larger treatment effects than global correction suggests, aligning with known annotation difficulty gradients for these categories.

\subsection{Robustness to Subgroup Misspecification}\label{sec:robustness}

A practical concern is sensitivity to the choice of subgroup variable $S$.
We test EC-HTE under three misspecification scenarios using synthetic data ($K=4$, $n_{\text{expert}}=500$): (a)~noisy subgroup assignment with 10--30\% label noise on $S$, (b)~coarsened subgroups ($K=4 \to K=2$ by merging heterogeneous pairs), and (c)~irrelevant subgroup variable (uncorrelated with misclassification heterogeneity).

\begin{table}[t]
\centering
\small
\begin{tabular}{lcccl}
\toprule
\textbf{Scenario} & \textbf{RMSE} & \textbf{Bias} & \textbf{Beats} & \textbf{Beats} \\
 & \textbf{Ratio} & \textbf{Ratio} & \textbf{Global?} & \textbf{Naive?} \\
\midrule
Correct spec. & 1.00 & 1.00 & yes & yes \\
Noisy 10\% & 1.08 & 1.12 & yes & yes \\
Noisy 20\% & 1.18 & 1.30 & marginal & yes \\
Noisy 30\% & 1.35 & 1.55 & no & yes \\
Coarsened $K\!=\!2$ & 1.45 & 1.40 & yes & yes \\
Irrelevant $S$ & 1.10 & 1.02 & $\approx$equal & yes \\
\bottomrule
\end{tabular}
\caption{Robustness analysis. RMSE and Bias Ratios are relative to correctly-specified EC-HTE. EC-HTE degrades gracefully: up to 20\% label noise retains advantage over global correction; irrelevant subgroups cause minimal harm due to HB shrinkage.}
\label{tab:robustness}
\end{table}

\paragraph{Graceful Degradation via Partial Pooling.}
Table~\ref{tab:robustness} shows that EC-HTE degrades smoothly under misspecification.
With 10\% subgroup label noise, RMSE increases by only 8\% and EC-HTE retains its advantage over global correction.
At 20\%, EC-HTE is marginally better.
The crossover occurs at $\sim$25--30\% noise, consistent with James-Stein theory: when the signal-to-noise ratio of subgroup differences drops below $\sim$1, pooling dominates stratification.

When the subgroup variable is irrelevant (uncorrelated with misclassification), HB partial pooling assigns $>$90\% shrinkage weight to the global mean (Table~\ref{tab:shrinkage}), making EC-HTE approximately equivalent to global correction with minimal overhead (5--15\% RMSE increase).
This is a key practical property: EC-HTE is safe to apply even when the practitioner is uncertain about subgroup relevance.

\begin{table}[t]
\centering
\small
\begin{tabular}{lccc}
\toprule
\textbf{Scenario} & \textbf{Mean} & \textbf{Min} & \textbf{Max} \\
 & \textbf{Shrinkage} & & \\
\midrule
Correct specification & 0.35 & 0.15 & 0.55 \\
Noisy 20\% & 0.55 & 0.35 & 0.70 \\
Noisy 30\% & 0.68 & 0.50 & 0.82 \\
Irrelevant subgroup & 0.92 & 0.85 & 0.97 \\
\bottomrule
\end{tabular}
\caption{HB shrinkage weights toward the global mean under misspecification. Higher values indicate more pooling. Under irrelevant subgroups, the model correctly pools almost entirely, collapsing to near-global correction.}
\label{tab:shrinkage}
\end{table}

Coarsening $K=4$ to $K=2$ retains partial heterogeneity structure: EC-HTE still beats naive correction, but the margin over global correction is halved, with 30--50\% of bias remaining uncorrected.
This motivates using fine-grained subgroups when available.

\subsection{Adaptive Expert Annotation Allocation}\label{sec:allocation}

When expert annotation budgets are fixed, allocation across subgroups matters.
We compare three strategies on synthetic data ($K=4$): (a)~uniform allocation (25\% per subgroup), (b)~uncertainty-weighted (proportional to posterior variance of confusion matrix estimates), and (c)~CATE-weighted (proportional to $|\hat{\tau}_k| \times$ misclassification uncertainty).

\begin{table}[t]
\centering
\small
\begin{tabular}{clccl}
\toprule
$n_{\text{exp}}$ & \textbf{Strategy} & \textbf{RMSE} & \textbf{Worst-SG} & \textbf{Alloc.} \\
 & & & \textbf{RMSE} & ($S_0$/$S_1$/$S_2$/$S_3$) \\
\midrule
\multirow{3}{*}{250} & Uniform & .385 & .62 & 25/25/25/25 \\
 & Uncertainty & .358 & .54 & 18/22/28/32 \\
 & CATE-wt. & .340 & .48 & 15/20/27/38 \\
\midrule
500 & Uniform & .245 & .38 & 25/25/25/25 \\
500 & CATE-wt. & .225 & .31 & 16/21/26/37 \\
\midrule
1000 & Uniform & .165 & .22 & 25/25/25/25 \\
1000 & CATE-wt. & .158 & .20 & 18/22/27/33 \\
\bottomrule
\end{tabular}
\caption{Expert annotation allocation strategies ($K\!=\!4$ extreme regime). CATE-weighted allocation concentrates labels on high-error minority subgroups ($S_2$, $S_3$), reducing worst-subgroup RMSE by 15--25\% at $n_{\text{expert}}\!=\!250$.}
\label{tab:allocation}
\end{table}

CATE-weighted allocation provides 8--15\% average RMSE reduction at $n_{\text{expert}}=250$ and 15--25\% worst-subgroup improvement (Table~\ref{tab:allocation}), with diminishing returns at larger budgets.
The strategy concentrates 40--60\% of the budget on the two highest-error subgroups ($S_2$, $S_3$), which comprise only 24\% of the data.
For the $K=4$ extreme regime, CATE-weighted allocation moves the crossover point (where EC-HTE RMSE matches global correction) from $n_{\text{expert}} \approx 1000$ to $\approx$750--800 --- helpful but not transformative.

\subsection{When to Use EC-HTE}\label{sec:when}

EC-HTE's advantage is conditional, not universal.
Based on our experiments, we provide per-$K$ budget recommendations:

\begin{enumerate}[leftmargin=*,itemsep=2pt]
\item \textbf{Audit first}: Estimate per-subgroup misclassification rates on a pilot sample. If $\Delta\pi < 0.05$, global correction is likely sufficient.
\item \textbf{$K=2$ subgroups}: EC-HTE is beneficial with $n_{\text{expert}} \geq 250$ when $\Delta\pi > 0.05$. Bias amplification from global correction is $\sim$1.3--1.7$\times$.
\item \textbf{$K=4$ subgroups}: EC-HTE requires $n_{\text{expert}} \geq 500$ (ensuring $\geq$60 labels for 12\% minority subgroups). With fewer labels, per-subgroup CM estimation noise exceeds the benefit of stratification. Bias amplification from global correction can reach 1.9--2.2$\times$ with heterogeneous LLMs.
\item \textbf{Low heterogeneity + small budget}: When $\Delta\pi < 0.05$ and $n_{\text{expert}} < 250$, global correction outperforms EC-HTE on RMSE. The diagnostic audit protocol identifies this scenario.
\end{enumerate}

\section{Conclusion}

We identify and formalize a bias mechanism in heterogeneous treatment effect estimation with LLM-annotated effect modifiers: global misclassification correction introduces bias proportional to $\Delta\pi \times \Delta\tau$, amplifying subgroup-level CATE errors by 1.3--2$\times$ with current SOTA LLMs.
EC-HTE addresses this through hierarchical Bayesian subgroup-stratified correction with partial pooling.
The framework degrades gracefully under subgroup misspecification, complements existing prediction-powered inference methods, and provides practitioners with auditing and budget allocation tools.

We emphasize that the moderate heterogeneity observed with current frontier LLMs on standard CSS tasks represents a lower bound: as LLM annotation expands to harder tasks --- cross-lingual hate speech detection, implicit framing analysis, nuanced political discourse coding --- misclassification heterogeneity will increase, making stratified correction increasingly consequential.
EC-HTE provides the methodological foundation for this transition.

\paragraph{Limitations.}
EC-HTE requires expert labels for confusion matrix estimation, adding annotation cost.
The RMSE tradeoff at small budgets ($n_{\text{expert}} < 250$ for $K=4$) means EC-HTE is not universally preferable.
Our real-data case study demonstrates pattern consistency but cannot validate against unknown ground-truth CATEs.
The $\Delta\pi \times \Delta\tau$ diagnostic requires a pilot labeling phase, though this cost is modest relative to the full annotation budget.

\bibliographystyle{plainnat}
\bibliography{references}

\appendix

\section{Hyperprior Sensitivity}\label{app:hyperprior}

We use $\text{Gamma}(2, 1)$ for the Dirichlet concentration parameter $\alpha_0$.
Sensitivity analysis with $2\times$ variations ($\text{Gamma}(4, 2)$ and $\text{Gamma}(1, 0.5)$) shows $<$5\% change in RMSE and $<$3\% change in bias across all experimental configurations, confirming robustness to hyperprior specification.

\section{Full Per-Subgroup Results for Multiple LLM Experiment}\label{app:llm_full}

Table~\ref{tab:per_subgroup_misclass} in the main text provides per-subgroup misclassification rates.
The rank order of heterogeneity across LLMs is Qwen3 $>$ GPT-4.1 $\geq$ Claude Sonnet~4.6.
Open-source models show higher heterogeneity and correspondingly greater EC-HTE benefit, likely due to less calibrated instruction-following across content types.
The correlation between $\Delta\pi$ and percentage bias reduction ($r = 0.85$--$0.95$) supports the $\Delta\pi \times \Delta\tau$ identity as a practical predictor of EC-HTE utility.

\section{Adaptive Allocation: Extended Results}\label{app:allocation}

At $n_{\text{expert}} = 1000$, all allocation strategies converge (Table~\ref{tab:allocation}): per-subgroup confusion matrices are well-estimated under any allocation.
The modest overall advantage of CATE-weighted allocation (2--5\% at $n=1000$) reflects that the fundamental bottleneck is total budget, not allocation efficiency.
Adaptive allocation adds $<$5\% computational overhead (one additional posterior inference step for allocation weights).

\end{document}
